% ============================================================
%  Postgraduate Dissertation -- ICML / AI Paper Style
%  Ockham's Causal Razor: Symbolic Regression for Dictionary
%  Discovery in Causal Inference
% ============================================================
\documentclass[11pt, a4paper]{article}

% --- Core Packages ---
\usepackage[margin=1.2in]{geometry}
\usepackage{times}
\usepackage{microtype}
\usepackage{parskip}
\usepackage{setspace}
\onehalfspacing

% --- Maths ---
\usepackage{amsmath, amssymb, amsthm, mathtools}
\usepackage{bm}

% --- Theorem Environments ---
\theoremstyle{plain}
\newtheorem{theorem}{Theorem}[section]
\newtheorem{proposition}[theorem]{Proposition}
\newtheorem{lemma}[theorem]{Lemma}

\theoremstyle{definition}
\newtheorem{definition}[theorem]{Definition}
\newtheorem{assumption}[theorem]{Assumption}

\theoremstyle{remark}
\newtheorem{remark}[theorem]{Remark}

% --- Figures, Tables, Highlighting ---
\usepackage{graphicx}
\usepackage{booktabs}
\usepackage{float}
\usepackage{caption}
\captionsetup{font=small, labelfont=bf}
\usepackage{subcaption}
\usepackage{multirow}
\usepackage{array}
\usepackage{colortbl}
\definecolor{highlight}{rgb}{1, 0.93, 0.62}
\newcommand{\hl}[1]{\cellcolor{highlight}#1}

% --- Algorithms ---
\usepackage{algorithm}
\usepackage{algpseudocode}

% --- Colours and Hyperlinks ---
\usepackage[dvipsnames]{xcolor}
\usepackage[colorlinks=true,
            linkcolor=MidnightBlue,
            citecolor=MidnightBlue,
            urlcolor=MidnightBlue]{hyperref}

% --- Bibliography ---
\usepackage[authoryear, round]{natbib}
\bibliographystyle{plainnat}

% --- Section Styling (ICML-like) ---
\usepackage{titlesec}
\titleformat{\section}{\large\bfseries\scshape}{\thesection.}{0.5em}{}
\titleformat{\subsection}{\normalsize\bfseries}{\thesubsection.}{0.5em}{}
\titleformat{\subsubsection}{\normalsize\itshape}{\thesubsubsection.}{0.5em}{}

% --- Header and Footer ---
\usepackage{fancyhdr}
\pagestyle{fancy}
\fancyhf{}
\renewcommand{\headrulewidth}{0.4pt}
\fancyhead[L]{\small\scshape Przemek Szkodon}
\fancyhead[R]{\small\scshape Ockham's Causal Razor}
\fancyfoot[C]{\thepage}

% --- Custom Commands ---
\newcommand{\lasso}{\textsc{Lasso}}
\newcommand{\pds}{\textsc{PDS}}
\newcommand{\dml}{\textsc{DML}}
\newcommand{\pysr}{\textsc{PySR}}
\newcommand{\epds}{\textsc{SR-PDS}}
\newcommand{\E}{\mathbb{E}}
\newcommand{\R}{\mathbb{R}}
\newcommand{\PP}{\mathbb{P}}
\newcommand{\norm}[1]{\left\lVert #1 \right\rVert}
\newcommand{\given}{\,|\,}
\newcommand{\Shat}{\hat{S}}
\newcommand{\betahat}{\hat{\beta}_0}
\newcommand{\Xstar}{\bm{x}^{*}}
\newcommand{\Hcal}{\mathcal{H}_{\mathcal{O}, C}}

% ============================================================
\begin{document}

\begin{center}
    {\Large\bfseries Ockham's Causal Razor:\\[0.3em]
    Symbolic Regression for Dictionary Discovery\\
    in Causal Inference}\\[1.2em]
    {\large Przemek Szkodon}\\[0.3em]
    {\normalsize Department of Economics and Data Science}\\
    {\normalsize University of Cambridge}\\[0.3em]
    {\normalsize \texttt{ps666@cam.ac.uk}}\\[1em]
    {\small \today}
\end{center}

\vspace{1em}

% ============================================================
\begin{abstract}
\noindent
This dissertation proposes a framework for using symbolic
regression in causal inference. The pattern, established in the
physical sciences \citep{schmidt2009gp, cranmer2020}, is to
let complex tools find structure and simple tools deliver
inference. We argue the same arrangement addresses a persistent
failure in high-dimensional causal estimation that the standard
treatment obscures: when the treatment depends nonlinearly on
covariates that also affect the outcome, the standard methods
carry bias on the order of the confounding magnitude across
$n=30$ to $n=5000$, with confidence intervals that contain the
true effect zero percent of the time. This holds for
Post-Double Selection with \lasso\ \citep{bch2014} and for
Double Machine Learning with either \lasso\ or neural-network
nuisances \citep{chernozhukov2018}. Our worked example is
\epds, which runs symbolic regression upstream of the \pds\
selection step. The justification is that symbolic regression
searches a finite class with bounded log-cardinality, delivering
parametric convergence rates where the neural-network analogue
scales quadratically in parameter count and does not bind at
applied sample sizes. \epds\ achieves coverage above $0.90$
where the standard methods collapse, whilst preserving the
named-confounder structure on which applied work depends.
\end{abstract}

\vspace{1em}

% ============================================================
\section{Introduction}
\label{sec:intro}

This dissertation argues for a particular relationship between
flexible machine learning and causal inference. The dominant
practice plugs flexible learners directly into causal procedures
and lets the asymptotic theory carry the inference. We argue
this is the wrong arrangement. Flexible methods belong at the
discovery step, where their representational power finds the
structure in the data. Simple, interpretable methods belong at
the inference step, where their tractability delivers standard
errors that an applied researcher can defend in peer review or
before a policy audience. The pattern has been productive in
the physical sciences for over a decade
\citep{schmidt2009gp, cranmer2020}, and we propose that empirical
economics has reason to follow suit. Our framework instantiates
the idea in the post-double-selection setting of
\citet{bch2014}, and the instance, which we call \epds, serves
as the worked example throughout the dissertation. The principle
behind it is more general; Section \ref{sec:discussion} sketches
extensions to matching, difference-in-differences, and
instrumental variables.

The case for the framework rests on a structural failure in
existing high-dimensional causal methods that the standard
treatment obscures. Applied causal work typically opens with a
long list of candidate controls and no clear theory for which
ones matter. Choosing among them is itself a statistical
problem, and the standard answer over the past decade has been
to lean on high-dimensional regularisation. The canonical
procedure is Post-Double Selection by \citet{bch2014}, which
runs two \lasso\ regressions, one for the outcome and one for
the treatment, takes the union of the selected variables, and
runs OLS on that union. Provided the nuisance functions are
approximately linear in the candidate covariates, the resulting
estimator of the treatment effect is asymptotically normal and
its confidence intervals have the right coverage.

That linearity assumption is the catch. If the propensity score
takes the form $m(\bm{x}) = 2 x_4 x_5$ rather than $m(\bm{x}) = 2
x_4 + 2 x_5$, the marginal correlation between $d$ and either $x_4$
or $x_5$ alone is essentially zero, and \lasso\ cannot find either
one. The same applies to any nonlinear functional form not spanned
by the original covariates, whether quadratics, products, or
transformations. The standard remedy is to expand the dictionary manually, by
including squares and selected interactions, but this makes the
analyst do the work the method was meant to do automatically, and
offers no guarantee that the chosen expansion contains the right
terms. Worse, the choice of expansion is itself a selection problem,
and a researcher who chooses badly inherits the omitted-variable
bias of the original procedure with none of the guarantees that
motivated using a regularised method in the first place.

Double Machine Learning by \citet{chernozhukov2018} provides a
different remedy. Rather than select variables, it replaces the
linear nuisance regressions with flexible machine-learning models,
takes orthogonal residuals, and regresses outcome residuals on
treatment residuals. In principle any sufficiently flexible learner
works, including neural networks. In practice, however, achieving
the required convergence rate with a neural network on a few
hundred observations is difficult, because the network has
thousands of parameters fit to a small training set, and early
stopping kicks in before the network has learned the relevant
nonlinear structure.
On the nonlinear-confounding designs we construct, \dml-NN's bias
is indistinguishable from naive OLS even at $n=5000$. \dml-\lasso\
inherits the same limitation as \pds-\lasso\ because the nuisance
estimator is still linear.

The two responses share a second, less remarked-upon weakness.
\dml\ produces nuisance estimates as black-box functions, leaving
a researcher applying the method with no list of confounders to
point at, to defend to a referee, or to compare against the next
study's choices. \pds\ preserves a named confounder set by construction,
which is part of why it remains the dominant procedure in applied
microeconometrics despite \dml's theoretical superiority. We are
reluctant to give up that property. The applied virtue of \pds\ is
not merely cosmetic. In peer review, in policy submissions, and in
subsequent meta-analyses, the ability to point at a specific list
of controls and defend it is a material feature of a method, not a
stylistic one.

The discovery-versus-inference distinction we propose has
precedent. \citet{schmidt2009gp} and \citet{cranmer2020} show
that complex machine-learning tools can be used to recover
simple structural laws from data: a flexible learner finds the
structure, that structure is represented as a compact symbolic
expression, and the symbolic expression carries the downstream
scientific work. Importing the pattern to causal econometrics
means that if symbolic regression can recover the functional
form of a nonlinear confounder, the recovered expression
belongs in a procedure that delivers valid inference, not in
the wings whilst a black-box learner produces the final
estimate. The simple representation is what enables the
scientific virtues of falsifiability, reproducibility, and
narrative that empirical economics has historically rested on.

We propose a procedure that keeps the \pds\ structure and adds the
flexibility \pds-\lasso\ lacks. The method, which we call \epds\
(symbolic-regression post-double selection), runs symbolic
regression on the
outcome and propensity equations to discover nonlinear basis terms,
adds those terms to the feature dictionary, and then runs standard
\pds-\lasso\ on the enriched dictionary. We use the \pysr\
implementation of symbolic regression \citep{cranmer2023pysr},
which performs evolutionary search over expression trees and
returns a Pareto frontier of expressions trading off complexity
against fit. The final estimator is OLS on a named set of selected
features, so the commitment-based interpretability of \pds\ is
preserved.

To evaluate the proposed method we construct eight data-generating
processes that span the regimes of interest: linear, quadratic,
interaction, mixed nonlinear, linear-confounded nonlinear, two
designs in which the propensity score depends nonlinearly on
covariates that also enter the outcome equation (mild and severe
nonlinear confounding), and a weak-signal design with mixed
continuous and binary covariates. We compare six methods (Full
OLS, \pds-\lasso, \dml-\lasso, \dml-RF, \dml-NN, \epds) on bias,
RMSE, coverage, and size, at fixed $n=500$ and across a
trajectory from $n=30$ to $n=50{,}000$. A cross-fitted variant of
\epds\ is included as a robustness check on the inferential
validity of the same-data procedure.

The empirical findings divide cleanly into three regimes. On
linear and lightly-nonlinear designs, \epds\ matches the
standard methods within Monte Carlo error: the flexibility
costs nothing when it is not needed, which is the right
baseline property for a proposed extension. On the
nonlinear-confounding designs the standard methods collapse
entirely, with bias exceeding $1.2$ on the mild design and
$1.5$ on the severe design and coverage of exactly zero;
\epds\ reduces bias by more than ninety-five
percent on the mild design and by more than eighty-five percent
on the severe design, with coverage at $0.9$ or above. The
trajectory analysis from $n=30$ to $n=5000$ shows that \dml's
failure on these designs is structural rather than small-sample:
\dml-NN's bias on the mild nonlinear-confounding design remains
above $1.0$ at $n=5000$, with no visible convergence. Section
\ref{sec:method} makes this precise. The convergence-rate
condition in \citet{chernozhukov2018} requires sample complexity
that scales with the squared parameter count of the nuisance
learner, which for a typical multilayer perceptron is several
orders of magnitude above what applied work delivers.

The dissertation is organised as follows. Section \ref{sec:lit}
positions the work within the related literature. Section
\ref{sec:background} reviews \pds, \dml, and symbolic
regression. Section \ref{sec:method} describes the \epds\
pipeline and the theoretical reasoning that motivates it. Section
\ref{sec:sim} describes the simulation design. Section
\ref{sec:results} presents results at fixed $n=500$ and across the
$n$-trajectory, with explicit reference back to the theoretical
predictions. Section \ref{sec:discussion} discusses limitations and
practical guidance. Section \ref{sec:conclusion} concludes.

% ============================================================
\section{Related Literature}
\label{sec:lit}

The literature on variable selection for causal inference begins
with the \lasso\ \citep{tibshirani1996}, and the recognition that
using it directly for inference is subtle, since the
regularisation that makes \lasso\ useful as a predictive method
induces an omitted-variable bias in any target parameter \citep{bch2014}. The post-%
double-selection procedure addresses this by running \lasso\
separately on the outcome and the treatment equations and taking
the union, which protects against omitted confounders missed by
either single regression. \citet{chernozhukov2018} generalise this
insight to a broad class of nuisance estimators via Neyman
orthogonality and cross-fitting, producing the Double Machine
Learning framework. Both procedures admit valid inference when
their assumptions hold, and both rest on sparsity assumptions
about the underlying nuisance functions, with the difference
being whether the sparsity is in a linear basis (\pds-\lasso) or
in whatever class the nuisance estimator can represent (\dml).

A parallel literature develops feature-selection methods for the
nonlinear case. LassoNet \citep{lemhadri2021lassonet} imposes a
hierarchical sparsity constraint linking a linear skip path to a
deep network. Deep Feature Selection \citep{chen2021dfs} uses an
$L_1$-penalised gating layer between the inputs and the main
network. Both rely on gradient descent and inherit the associated
sensitivities to architecture, initialisation, and local optima.

Symbolic regression itself dates to \citet{koza1992}'s genetic
programming and has been brought into modern data analysis by
\citet{cranmer2020}, who use a graph neural network to compress
relational data before symbolic search, and the \pysr\
implementation \citep{cranmer2023pysr}, which provides a high-%
performance engine for evolutionary search over expression trees.
The key practical fact for our purposes is that symbolic
regression runtime scales rapidly with input dimensionality. With
$p=50$ candidate covariates, a single \pysr\ fit takes
approximately twenty seconds. The maximum-complexity budget and
the parsimony coefficient control the search and limit spurious
low-coefficient terms.

The validity of \pds-\lasso\ inference rests on the basis
dictionary being fixed in advance, not estimated from the data
used for the final OLS step. Our pipeline violates this
assumption. \citet{lee2016exact} and \citet{tibshirani2018uniform}
develop exact and uniform post-selection inference for the
\lasso, but extending these to a symbolic-regression-discovered
dictionary requires care. We discuss the implications and offer
simulation evidence on the empirical coverage in Section
\ref{sec:results}, and we leave the formal extension to future
work. The natural fix is sample splitting, which restores formal
validity at the cost of statistical efficiency. We have not
implemented sample splitting in the current pipeline, which is
the principal outstanding inferential question. We discuss this
further in Section \ref{sec:discussion}.

Two recent simulation studies are particularly close to our
setting. \citet{fuhr2024} compare \dml\ with \lasso, random
forest, and neural-network nuisance estimators across a range
of nonlinear designs, and document bias on the order of $0.1$
to $0.3$ at $n=2000$ with neural-network nuisances. Their
designs do not include nonlinear confounding of the kind we
construct in our DGPs 6 and 7, where the propensity and outcome
share nonlinear terms; the corresponding bias on our designs is
roughly five times larger. \citet{klusowski2024doubledeep}
examine \dml\ with deep nuisance estimators and identify the
convergence-rate condition as the binding constraint at applied
sample sizes, providing the formal theory that motivates our
sample-complexity calculation in Section \ref{sec:theory}. Our
results extend both papers in two directions. We construct
designs where the standard methods fail completely rather than
mildly, and we trace the failure across $n=30$ to $n=5000$ to
show that it does not resolve with more data. Neither paper
proposes a fix.

The closest methodological neighbours are post-selection
inference and feature selection for nonlinear models.
\citet{lee2016exact} and \citet{tibshirani2018uniform} develop
post-selection inference machinery that could in principle be
adapted to a data-estimated dictionary, though neither paper
considers symbolic regression as the selection mechanism.
\citet{lemhadri2021lassonet} and \citet{chen2021dfs} provide
neural-network analogues of \lasso\ for nonlinear feature
selection; both inherit the optimisation pathologies of
gradient-trained networks and produce black-box estimates
without the named-confounder property that our procedure
preserves. The Stein-based selector of \citet{duetal2025} is
gradient-free and deterministic but does not perform
functional-form recovery; it identifies which variables matter
nonlinearly but not how. Our framework combines its strengths
as an upstream filter with the symbolic recovery that follows.

The gap we fill is the absence of an existing method that
performs functional-form recovery for nuisance estimation within
the post-double-selection inference structure of \pds-\lasso.
Adjacent tools such as the sensitivity analysis of
\citet{cinelli2020sensitivity} and the partial identification
of \citet{conley2012plausibly} are complementary; they assess
the robustness of causal estimates to assumed violations of
identifying assumptions, whilst our procedure reduces the gap
between the analyst's specification and the data-generating
process before such sensitivity analysis becomes necessary.

% ============================================================
\section{Background}
\label{sec:background}

We work in the partially linear model of \citet{robinson1988},
\begin{equation}
    y_i = \beta_0 d_i + g(\bm{x}_i) + \varepsilon_i, \qquad
    d_i = m(\bm{x}_i) + v_i,
    \label{eq:plm}
\end{equation}
where $y_i$ is the outcome, $d_i$ the treatment, $\bm{x}_i \in
\R^p$ the candidate controls, and $g, m$ unknown nuisance
functions. We observe iid draws of $(y_i, d_i, \bm{x}_i)$ and
seek to estimate $\beta_0$.

\subsection{Post-Double Selection \lasso}

The \pds-\lasso\ procedure runs $\lasso(y \sim \bm{x})$ to obtain
a selected set $\Shat_y$, runs $\lasso(d \sim \bm{x})$ to obtain
$\Shat_d$, and regresses $y$ on $d$ and $\bm{x}_{\Shat_y \cup
\Shat_d}$ by OLS. The regularisation parameter $\lambda$ is set
to the theoretical bound $\lambda \geq c \sigma n^{-1/2}
\Phi^{-1}(1 - \alpha/(2p))$ for some constant $c > 1$, which
ensures the false-discovery rate is controlled. \citet{bch2014}
prove that under approximate sparsity of $g$ and $m$, the
resulting $\betahat$ is asymptotically normal at the parametric
rate with valid standard errors. The procedure preserves a named
control set, an important feature for applied work.

The assumption to watch is approximate sparsity in $\bm{x}$. If
the true $g$ contains a term $5 x_1^2 x_2$, no element of
$\bm{x}$ captures this directly, and \lasso\ on the raw $\bm{x}$
will fail to select $\{x_1, x_2\}$ on grounds of their marginal
correlation with $y$. The omission propagates to $\betahat$ as
confounding bias. The theoretical structure of the \pds\ result
is robust to a small amount of nonlinearity (the ``approximate''
in approximate sparsity permits a vanishing approximation error),
but the empirical behaviour of the procedure under substantial
nonlinear confounding is empirically comparable to that of
naive OLS.

\subsection{Double Machine Learning}

\dml\ replaces the \lasso\ selection step with two flexible
nuisance estimators $\hat g$ and $\hat m$, constructs residuals
$\tilde y = y - \hat g(\bm{x})$ and $\tilde d = d - \hat
m(\bm{x})$, and estimates $\beta_0$ as the OLS coefficient of
$\tilde y$ on $\tilde d$. Cross-fitting on disjoint folds removes
the overfitting concern. The estimator is asymptotically normal
provided the nuisance estimators converge at rate $\norm{\hat g -
g}_2 \cdot \norm{\hat m - m}_2 = o_p(n^{-1/2})$. This product-%
rate condition accommodates a wide range of learners in
principle. In practice, at the sample sizes typical of applied
microeconometric work, a multilayer perceptron with several
thousand parameters fit to a few hundred observations does not
approach the required rate for genuinely nonlinear $g$ and $m$,
and the estimator inherits the bias. The trajectory results in
Section \ref{sec:results} provide direct empirical evidence on
this point.

\subsection{Symbolic Regression}

\pysr\ \citep{cranmer2023pysr} performs evolutionary search over
expression trees with user-specified binary and unary operators
(we use $\{+, -, \cdot, \log, \sqrt{\cdot}\}$), a maximum
complexity budget, and a parsimony coefficient that penalises
complex expressions. The output is a Pareto frontier of
expressions trading off complexity against fit, and we select a
single expression by minimising $\ell + \rho c$, where $\ell$ is
validation loss and $c$ is complexity. The parsimony coefficient
$\rho$ controls the trade-off and must be tuned to the
application. We use $\rho = 0.005$ for outcome fits and $\rho =
0.01$ for propensity fits, reflecting the typically stronger
signal in the outcome equation.

The key practical concern at large $p$ is the size of the
expression-tree search space. With fifty candidate variables,
\pysr\ can select low-coefficient spurious terms involving
inactive features; the maximum-complexity budget and the
parsimony coefficient are the two levers that control this.

% ============================================================
\section{Methodology}
\label{sec:method}

The \epds\ pipeline runs symbolic regression on the outcome and
the propensity equations, extracts the discovered nonlinear basis
terms, appends them to the original feature matrix, and runs
standard \pds-\lasso\ on the enriched dictionary. Algorithm
\ref{alg:epds} states the procedure.

\begin{algorithm}[H]
\caption{Symbolic-Regression Post-Double Selection (\epds)}
\label{alg:epds}
\begin{algorithmic}[1]
    \Require Data $(y, d, \bm{x})$ with $\bm{x} \in \R^{n \times p}$
    \State Fit \pysr$(y \sim \bm{x})$ and \pysr$(d \sim \bm{x})$;
           let $\mathcal{T}_y$ and $\mathcal{T}_d$ be the sets of
           extracted nonlinear terms (e.g.\ $x_i^2$, $x_i x_j$,
           $\log|x_i|$).
    \State Construct $\Xstar = [\bm{x},\
           \phi_{\mathcal{T}_y \cup \mathcal{T}_d}(\bm{x})]$, the
           original $\bm{x}$ augmented with one numerical column
           per discovered term.
    \State Run $\lasso(y \sim \Xstar)$ and $\lasso(d \sim \Xstar)$
           with $\lambda^*$ scaled to $p^* = \dim\Xstar$. Let
           $\Shat^*$ be the union of the two selected sets.
    \State Compute $\betahat$ as the OLS coefficient on $d$ in
           the regression of $y$ on $d$ and $\Xstar_{\Shat^*}$,
           with HC1 standard errors.
\end{algorithmic}
\end{algorithm}

\subsection{Symbolic Regression Configuration}

We configure \pysr\ with operator set $\{+, -, \cdot, \log,
\sqrt{\cdot}\}$, maximum complexity 30, and 40 evolutionary
iterations. The operator set is deliberately small. We exclude
division because it produces unstable expressions that introduce
numerical noise into the downstream regressions, and we exclude
exponentiation because the values that the optimiser discovers as
exponents are not generally interpretable as economically
meaningful. The maximum complexity of 30 nodes comfortably
accommodates expressions of the form $a_1 x_1^2 x_2 + a_2 \log|x_3| + a_3 x_4
x_5$ (complexity 13) plus deeper compositions that may arise in
applied settings, whilst the parsimony coefficient keeps the
chosen expression near the elbow of the Pareto frontier.

The parsimony coefficient is the most consequential
hyperparameter. We set $\rho_y = 0.005$ for the outcome equation
and $\rho_d = 0.01$ for the propensity equation. The asymmetry
reflects the different signal-to-noise ratios in the two
equations in our DGPs. The outcome signal is typically four to
five times larger in magnitude than the propensity signal, and a
parsimony coefficient that is appropriate for the outcome would
underfit the propensity. We arrived at these values through pilot
runs in which we observed the propensity equation occasionally
collapsing to a single-term expression at $\rho_d = 0.005$.

\subsection{Term Extraction}

The output of \pysr\ is a symbolic expression. To use it as a
basis we walk the expression tree and extract individual nonlinear
terms, classified into a small taxonomy: squares ($x_i^2$),
interactions ($x_i x_j$ with $i < j$), cubic-style products
($x_i^2 x_j$), logarithms ($\log|x_i|$), and square roots
($\sqrt{|x_i|}$). Pure monomials, polynomials of higher degree,
and operator combinations outside this taxonomy are not currently
extracted. In our DGPs this restriction does not bind, but we
mention it for completeness.

For the logarithmic terms specifically, we construct the column
$\log(|x_i| + 10^{-6})$ to avoid evaluating the logarithm at
zero. The discrepancy between this construction and \pysr's
symbolic $\log(x_i)$ is negligible for $|x_i| \gg 10^{-6}$ and
does not bind in practice. The issue that does bind is that
\pysr's \texttt{log} operator does not natively handle negative
arguments, so for a data-generating $\log|x_3|$ with $x_3 \sim
N(0, 1)$, the term must be discovered as $\log(x_3^2)$ rather
than directly. We document this issue in the limitations.

\subsection{The Enriched \lasso\ Step}

We run two \lasso\ regressions on $\Xstar$ using the theoretical
$\lambda^*$ scaled to the enriched dimension $p^*$:
\begin{equation}
    \lambda^* = \frac{c}{\sqrt{n}} \cdot
    \Phi^{-1}\!\left(1 - \frac{\alpha}{2 p^*}\right),
    \qquad c = 1.1, \quad \alpha = 0.05.
\end{equation}
The union of the two selected sets becomes the final control set,
and the OLS regression of $y$ on $d$ and that set delivers the
point estimate and confidence interval. The choice $c = 1.1$
matches the recommended value in \citet{bch2014}. Because $p^*$
exceeds $p$ by the number of discovered nonlinear terms, the
enriched $\lambda^*$ is fractionally larger than the original
\pds\ $\lambda$. In our DGPs with $p^* - p \in \{2, 3, 4\}$ this
fractional adjustment is below one percent.

\subsection{Theoretical Justification: Why \epds\ Works at Moderate $n$}
\label{sec:theory}

Why should \epds\ work in regimes where neural-network \dml\
does not? The intuition is one of search-space size. \pysr\
looks for the right answer within a small, structured space of
mathematical expressions, namely those of bounded complexity
built from a handful of operators. A neural network, by
contrast, searches a much larger space of functions representable
by its parameters. When the truth lives in a small space, an
estimator that knows where to look needs much less data to find
it than an estimator that does not. The argument has three
substantive steps and a synthesis.

\subsubsection{Finite-class ERM gives parametric rates}

Fix an
operator set $\mathcal{O}$ and a maximum complexity $C$, and let
$\Hcal$ denote the set of all expression trees of complexity at
most $C$ built from operators in $\mathcal{O}$ and the $p$
candidate variables. The set $\Hcal$ is finite, with cardinality
bounded above by $|\mathcal{O}|^C \cdot p^C$, so $\log |\Hcal| =
O(C \log(|\mathcal{O}| \cdot p))$. For our setting with
$|\mathcal{O}| = 5$, $C = 30$, and $p = 50$, this gives $\log
|\Hcal|$ on the order of two hundred. Once the search space is
finite and not too large, the classical sample-complexity
result for empirical risk minimisation applies.

\begin{samepage}
\begin{proposition}[Finite-class ERM rate]
\label{prop:erm}
Let $\Hcal$ be a finite hypothesis class, and suppose the true
nuisance function satisfies $g \in \Hcal$. The empirical risk
minimiser
$\hat{g} = \arg\min_{f \in \Hcal} \tfrac{1}{n}
\sum_{i=1}^n (y_i - f(\bm{x}_i))^2$
satisfies, with probability at least $1 - \delta$,
\begin{equation*}
    \norm{\hat{g} - g}^2_{L^2}
    \;\leq\;
    \frac{c \, \sigma^2 \log(|\Hcal|/\delta)}{n},
\end{equation*}
for an absolute constant $c$ depending on the sub-Gaussian
noise variance proxy $\sigma^2$.
\end{proposition}

\emph{Sketch.} For each $f \in \Hcal$ the empirical squared loss
concentrates around the population squared loss at rate
$1/\sqrt{n}$ by Hoeffding's inequality, and a union bound over
the $|\Hcal|$ candidate expressions yields the stated rate. See
\citet{wainwright2019high}, Chapter 4. $\square$
\end{samepage}

In plain language, when the relationship in the data can be
written as one of the finitely many expressions \pysr\
considers, the number of observations required to find it grows
only with the logarithm of how many such expressions there are.
The contrast with neural networks is sharp. For a multilayer
perceptron with parameter count $P$, the analogous guarantee
scales as $n = \Omega(P^2 \log P / \epsilon^2)$
\citep{anthony2009neural}. Our $(64, 32)$ MLP has $P \approx
5000$ parameters, so the implied $n$ for $\epsilon = 0.1$
accuracy is on the order of $10^6$, three orders of magnitude
above the sample sizes we consider. The same calculation in our
finite-class setting gives an implied $n$ on the order of
$10^3$, comfortably within applied range. The trajectory result
in Section \ref{sec:results} that \dml-NN bias does not decline
with $n$ across our range is the empirical signature of this
mismatch.

The proposition assumes the truth lies in $\Hcal$.
Misspecification is the more common applied case, and the
following remark records what happens then.

\begin{remark}[Misspecified ERM rate]
\label{rem:misspec}
Suppose the true nuisance function $g \notin \Hcal$, with
approximation error
$\Delta^2 := \inf_{f \in \Hcal} \norm{f - g}^2_{L^2}$.
Then under the conditions of Proposition \ref{prop:erm}, the
ERM minimiser $\hat{g}$ satisfies
\begin{equation*}
    \norm{\hat{g} - g}^2_{L^2}
    \;\leq\;
    \Delta^2 \;+\; \frac{c \, \sigma^2 \log(|\Hcal|/\delta)}{n}
\end{equation*}
with probability at least $1-\delta$.
\end{remark}

The remark separates the error into two components. The
estimation error $\log|\Hcal|/n$ vanishes as $n$ grows. The
approximation error $\Delta^2$ does not. For \dml\ inference
to engage, the product-rate condition $\norm{\hat{g} - g}_2
\cdot \norm{\hat{m} - m}_2 = o_p(n^{-1/2})$ requires
$\Delta_g \cdot \Delta_m = o_p(n^{-1/2})$, which fails for any
fixed approximation error as $n$ grows. The misspecified case
therefore does not just slow convergence, it breaks it. DGP-4
inhabits this regime with respect to the $\log|x_3|$ term given
our \pysr\ operator-set configuration; the persistent bias floor
we observe is the approximation error rather than a slow rate
of convergence, and the only fix is to enlarge $\Hcal$ to cover
the truth.

\subsubsection{Genetic-programming search converges to the ERM}

\pysr\ does not enumerate $\Hcal$ exhaustively. It performs
evolutionary search through mutation, crossover, and tournament
selection. \citet{schmidt2009gp} establish that, under standard
conditions on population size and Pareto-based selection, the
probability of recovering the true expression converges to one
as the iteration budget grows, provided the truth lies in the
searched class. \citet{cranmer2023pysr} document this
empirically across a broad range of physical-law discovery
problems.

Proposition \ref{prop:erm} therefore applies to \pysr's output
only when the iteration budget is large enough for the search
to succeed. In our DGPs, with 40 iterations and the default
population size, recovery is essentially one for low-complexity
targets such as products and squares, but drops below one for
high-complexity targets such as the cubic $x_1^2 x_2$ in
DGP-7. This is visible in the term-recovery panels in Section
\ref{sec:results} and explains the slower convergence of \epds\
bias on DGP-7 relative to DGP-6.

\subsubsection{The Pareto elbow selects the correct expression}

\pysr\ returns the full Pareto frontier $\{(c, \ell_c)\}$
trading expression complexity $c$ against fit $\ell_c$.
Selecting the expression that minimises $\ell_c + \rho c$
corresponds to selecting the point at the elbow of the curve,
where marginal complexity reduction yields disproportionate fit
gain. The parsimony coefficient $\rho$ plays the role of a BIC
log-$n$ term, acting as the price the expression pays for each
additional node of complexity.

When the truth has complexity $c^*$ and lies in $\Hcal$, the
loss profile is piecewise. $\ell_c$ is roughly constant and
equal to the noise variance for $c \geq c^*$, and sharply
higher for $c < c^*$ where the expression cannot represent the
truth. The elbow is at $c^*$, and any reasonable $\rho$ that
straddles the kink will select the correct complexity. When the
truth lies outside $\Hcal$, the loss profile decreases
gradually with $c$ and the elbow is less sharply identified,
making the parsimony coefficient more consequential. This is a
second mechanism by which DGP-4 underperforms relative to DGPs
6 and 7.

\subsubsection{Synthesis}

On a DGP where the
true nuisance functions lie in $\Hcal$ at low complexity and
the \pysr\ search budget is adequate, the discovered
expression coincides with the truth and the enriched
dictionary contains a basis that spans the relevant nonlinear
structure. The downstream \pds-\lasso\ step is then operating
on a dictionary that effectively includes the truth, and the
original \citet{bch2014} inference machinery applies. When the
truth lies outside $\Hcal$, the discovered expression is the
best approximation within the class, and the bias is
determined by the approximation error. When the truth lies in
$\Hcal$ at high complexity, the search budget binds and
convergence is gradual. The three regimes are visible in the
next section's results, and this story explains why each DGP
behaves as it does.

\subsection{Inferential Considerations}
\label{sec:inference}

The \citet{bch2014} inference theorem assumes the basis
dictionary is fixed in advance. We estimate it from the same
data we use for the final regression. Formally, this is a post-%
selection inference problem to which we do not yet have a
complete theoretical answer.

Three considerations bear on the practical inferential validity
of \epds. First, the dictionary expansion is small. In our DGPs,
\pysr\ typically discovers between two and five nonlinear terms,
so $\Xstar$ adds at most a few columns to the original feature
space. The flexibility added by symbolic regression is modest
relative to the flexibility already implicit in the \lasso\
search over the raw $\bm{x}$. Second, the parsimony coefficient
ensures the symbolic regression step is itself sparse, with
spurious terms of small coefficients penalised out of the
discovered expression, so the basis we add to the dictionary is
itself a sparse selection over expression space. Third, simulation
evidence in Section \ref{sec:results} shows empirical coverage
close to or above the nominal $95\%$ on the designs where \epds\
performs well, including the nonlinear-confounding designs where
the standard methods fail completely. We interpret this as
empirical evidence that the inferential cost of estimating the
basis is small. We do not claim that this evidence substitutes
for a formal proof. The principled extension would be a Belloni-%
style theorem accommodating data-estimated dictionaries, perhaps
via the uniform post-selection inference machinery of
\citet{tibshirani2018uniform} or via a sample-splitting argument.
We outline this programme in the conclusion and leave it to
future work.

% ============================================================
\section{Simulation Design}
\label{sec:sim}

\subsection{Data-Generating Processes}

We construct seven DGPs that progressively introduce nonlinearity
and nonlinear confounding. All use $n=500$, $p=50$, sparsity
$s=6$, and true treatment effect $\beta_0 = 0.5$. Covariates are
drawn iid as $\bm{x}_i \sim N(\bm{0}, I_{50})$, and noise terms
are iid standard normal. The outcome equation always takes the
form $y_i = \beta_0 d_i + g(\bm{x}_i) + \varepsilon_i$. Table
\ref{tab:dgps} summarises $g$ and the propensity $m$ for each
design.

\begin{table}[H]
\centering
\caption{The seven data-generating processes. Indices are 1-based.}
\label{tab:dgps}
\small
\begin{tabular}{lll}
    \toprule
    DGP & $g(\bm{x})$ & $m(\bm{x})$ \\
    \midrule
    1 Linear
        & $5 x_1 - 5 x_2 + 3 x_3 - 3 x_4 + x_5 - x_6$
        & $0$ \\
    2 Quadratic
        & $5 x_1^2 - 5 x_2 + 3 x_3^2 - 3 x_4 + x_5$
        & $0$ \\
    3 Interactions
        & $5 x_1 x_2 - 5 x_3 + 3 x_4 x_5 - 3 x_6$
        & $0$ \\
    4 Mixed NL
        & $5 x_1^2 x_2 - 5 \log|x_3| + 3 x_4 x_5 - 3 x_6$
        & $0$ \\
    5 NL + linear conf.
        & DGP-4 $g$
        & $x_1 + x_2 + x_3 + x_4 + x_5 + x_6$ \\
    6 NL + mild NL conf.
        & DGP-4 $g$
        & $2 x_4 x_5$ \\
    7 NL + severe NL conf.
        & DGP-4 $g$
        & $2 x_1^2 x_2 + 3 x_4 x_5$ \\
        8 Weak NL conf. (dummies)
        & {\small TBD} & {\small TBD} & {\small TBD} & {\small TBD} & {\small TBD} & {\small TBD} & {\small TBD} \\
    \bottomrule
\end{tabular}
\end{table}

DGPs 1 through 5 follow the standard simulation literature.
DGP-1 is the canonical \pds-\lasso\ benchmark. DGPs 2 and 3
introduce univariate nonlinearity and interactions. DGP-4 is
inspired by Case 5 of \citet{duetal2025}. DGP-5 layers a linear
propensity score over the mixed nonlinear outcome, and is the
standard DGP for testing \pds-style methods under confounding.

DGPs 6 and 7 are our contribution to the simulation design. The
propensity score depends on covariates that also appear in the
outcome equation, through a nonlinear functional form. This is
precisely the failure mode that motivates the development of
\epds. We can compute the analytical confounding strength for
each:
\begin{align}
    \text{DGP-6:}\quad
    & \text{Cov}(D, g(\bm{X})) = 2 \cdot 3 \cdot \text{Var}(x_4 x_5) = 6, \\
    \text{DGP-7:}\quad
    & \text{Cov}(D, g(\bm{X})) = 2 \cdot 5 \cdot \text{Var}(x_1^2 x_2) + 3 \cdot 3 \cdot \text{Var}(x_4 x_5) = 30 + 9 = 39.
\end{align}
The implied naive-OLS bias is $6/5 = 1.20$ for DGP-6 and
$39/22 \approx 1.77$ for DGP-7. The simulated values of
$1.244$ and $1.710$ reported in Section \ref{sec:results} match
these within Monte Carlo error.

The construction of DGPs 6 and 7 deserves explicit comment. We chose the shared nonlinear terms ($x_4 x_5$ in
DGP-6, plus $x_1^2 x_2$ and $x_4 x_5$ in DGP-7) so that the
confounding channel is precisely the kind of structure a linear
\lasso\ cannot capture. The shared term $x_4 x_5$ has zero
marginal correlation with either $x_4$ or $x_5$ in isolation, so
the propensity-equation \lasso\ cannot select these variables as
confounders. The same logic applies to $x_1^2 x_2$ in DGP-7. The
standard methods are in this sense designed against in the
propensity equation, not because we chose unrealistic functional
forms, but because we chose precisely the functional forms that
the propensity \lasso\ cannot represent. Remark
\ref{rem:misspec} formalises this: the truth lies outside the
linear hypothesis class of \pds-\lasso, so the ERM bound
carries an approximation error proportional to the confounding
magnitude that does not vanish with $n$.

\subsection{Methods Compared}

We compare seven methods:
\begin{enumerate}
    \item \textbf{Full OLS (infeasible)}: $y$ on $d$ and all $p$
          raw covariates. Singular when $p \geq n$, in which case
          it falls back to \pds-\lasso.
    \item \textbf{\pds-\lasso} \citep{bch2014}.
    \item \textbf{\dml-\lasso} \citep{chernozhukov2018}, cross-%
          fitted over five folds.
    \item \textbf{\dml-RF}: \dml\ with a random-forest nuisance
          estimator (500 trees, default depth, minimum leaf
          size 5).
    \item \textbf{\dml-NN}: \dml\ with a $(64, 32)$ MLP nuisance
          estimator, ReLU activation, early stopping on a $10\%$
          validation hold-out.
    \item \textbf{\epds}: Algorithm \ref{alg:epds}.
    \item \textbf{\epds\ (CF)}: cross-fitted \epds. For each of
          five folds, \pysr\ is fit on the training portion and
          \lasso\ predictions are made on the held-out fold;
          residuals are aggregated for the final OLS. Used as a
          robustness check on the inferential validity of the
          same-data variant.
\end{enumerate}

\subsection{Hyperparameters}

\pysr\ uses operator set $\{+, -, \cdot, \log, \sqrt{\cdot}\}$,
maximum complexity 30, parsimony 0.005 for outcome fits and 0.01
for propensity fits, 40 iterations, random seed 42.

The neural-network nuisance estimator for \dml-NN uses scikit-%
learn's \texttt{MLPRegressor} with the default optimizer (Adam),
hidden architecture $(64, 32)$, ReLU activation, learning rate
$0.001$, maximum 500 iterations, $10\%$ validation fraction for
early stopping, and patience of 10. These are essentially the
defaults, and we did not tune them per DGP. A reviewer may
reasonably ask whether better hyperparameters would rescue
\dml-NN on the nonlinear-confounding designs, and we discuss
this question in Section \ref{sec:discussion}.

\subsection{Evaluation}

We report bias $\E[\betahat - \beta_0]$, RMSE $\sqrt{\E[(\betahat
- \beta_0)^2]}$, coverage of nominal $95\%$ confidence intervals,
and size (rejection rate of $H_0: \beta_0 = 0$). We compute each
metric over $n_{\text{rep}}$ Monte Carlo replications.

\paragraph{Pilot replication count.} The results below use
$n_{\text{rep}} = 10$ for the fixed-$n$ design and $n_{\text{rep}}
= 3$ for the trajectory analysis. The Monte Carlo standard error
on the coverage estimate is approximately $\pm 0.16$ at
$n_{\text{rep}} = 10$ and wider at $n_{\text{rep}} = 3$. The headline effects we emphasise are larger than this margin by
an order of magnitude, and the finer comparisons should be
re-examined under the planned full production run at
$n_{\text{rep}} = 500$.

% ============================================================
\section{Results}
\label{sec:results}

\subsection{Fixed-$n$ Results}

Tables \ref{tab:bias} through \ref{tab:coverage} report bias, RMSE,
coverage, and size at $n=500$, across all eight DGPs and the
methods of Section \ref{sec:sim}. The cells highlighted in
yellow are the \epds\ values on the two nonlinear-confounding
DGPs. They are the headline result of this dissertation. In the
regime where all standard methods fail completely, \epds\ is the
only method that works.

\begin{table}[H]
\centering
\caption{Bias $\E[\betahat - \beta_0]$ at $n=500$. True $\beta_0
= 0.5$. Highlighted cells are the \epds\ values on
DGPs 6 and 7. Pilot results, $n_{\text{rep}} = 10$; the
production run regenerates these at $n_{\text{rep}} = 500$.}
\label{tab:bias}
\small
\setlength{\tabcolsep}{4pt}
\begin{tabular}{lrrrrrrr}
    \toprule
    DGP
        & Full
        & \pds-\lasso
        & \dml-\lasso
        & \dml-RF
        & \dml-NN
        & \epds
        & \epds (CF) \\
    \midrule
    1 Linear
        & $-0.007$ & $-0.019$ & {\small TBD} & {\small TBD} & $\phantom{-}0.054$ & $-0.007$ & {\small TBD} \\
    2 Quadratic
        & $-0.128$ & $-0.117$ & {\small TBD} & {\small TBD} & $-0.077$ & $-0.065$ & {\small TBD} \\
    3 Interactions
        & $\phantom{-}0.002$ & $\phantom{-}0.006$ & {\small TBD} & {\small TBD} & $\phantom{-}0.115$ & $\phantom{-}0.023$ & {\small TBD} \\
    4 Mixed NL
        & $\phantom{-}0.080$ & $\phantom{-}0.112$ & {\small TBD} & {\small TBD} & $\phantom{-}0.216$ & $\phantom{-}0.005$ & {\small TBD} \\
    5 NL + linear conf.
        & $\phantom{-}0.082$ & $\phantom{-}0.115$ & {\small TBD} & {\small TBD} & $\phantom{-}0.238$ & $\phantom{-}0.043$ & {\small TBD} \\
    6 NL + mild NL conf.
        & $\phantom{-}1.218$ & $\phantom{-}1.262$ & {\small TBD} & {\small TBD} & $\phantom{-}1.204$ & \hl{$-0.033$} & {\small TBD} \\
    7 NL + severe NL conf.
        & $\phantom{-}1.550$ & $\phantom{-}1.594$ & {\small TBD} & {\small TBD} & $\phantom{-}1.528$ & \hl{$\phantom{-}0.245$} & {\small TBD} \\
    8 Weak NL conf. (dummies)
        & {\small TBD} & {\small TBD} & {\small TBD} & {\small TBD} & {\small TBD} & {\small TBD} & {\small TBD} \\
    \bottomrule
\end{tabular}
\end{table}

\begin{table}[H]
\centering
\caption{RMSE of $\betahat$ at $n=500$. Highlighted cells are
the \epds\ values on DGPs 6 and 7. Pilot results,
$n_{\text{rep}} = 10$.}
\label{tab:rmse}
\small
\setlength{\tabcolsep}{4pt}
\begin{tabular}{lrrrrrrr}
    \toprule
    DGP
        & Full
        & \pds-\lasso
        & \dml-\lasso
        & \dml-RF
        & \dml-NN
        & \epds
        & \epds (CF) \\
    \midrule
    1 Linear
        & $0.035$ & $0.039$ & $0.050$ & {\small TBD} & $0.149$ & $0.040$ & {\small TBD} \\
    2 Quadratic
        & $0.227$ & $0.231$ & $0.251$ & {\small TBD} & $0.341$ & $0.135$ & {\small TBD} \\
    3 Interactions
        & $0.235$ & $0.234$ & $0.230$ & {\small TBD} & $0.359$ & $0.137$ & {\small TBD} \\
    4 Mixed NL
        & $0.331$ & $0.314$ & $0.336$ & {\small TBD} & $0.430$ & $0.357$ & {\small TBD} \\
    5 NL + linear conf.
        & $0.331$ & $0.317$ & $0.321$ & {\small TBD} & $0.397$ & $0.323$ & {\small TBD} \\
    6 NL + mild NL conf.
        & $1.248$ & $1.223$ & $1.269$ & {\small TBD} & $1.223$ & \hl{$0.244$} & {\small TBD} \\
    7 NL + severe NL conf.
        & $1.564$ & $1.554$ & $1.597$ & {\small TBD} & $1.531$ & \hl{$0.677$} & {\small TBD} \\
        8 Weak NL conf. (dummies)
        & {\small TBD} & {\small TBD} & {\small TBD} & {\small TBD} & {\small TBD} & {\small TBD} & {\small TBD} \\
    \bottomrule
\end{tabular}
\end{table}

The bias and RMSE results divide cleanly into the three regimes
predicted by the analysis in Section \ref{sec:theory}.

On the linear and lightly-nonlinear DGPs (1 through 4), all
methods perform acceptably, with the \epds\ variants matching or
slightly improving on \pds-\lasso. The improvement on the
interaction DGP is substantial: \epds\ achieves a small bias
and RMSE of $0.037$, compared with \pds-\lasso's $0.002$ and
$0.234$. The RMSE reduction reflects the narrower variance of the
enriched estimator. Once \pysr\ discovers the $x_1 x_2$ and
$x_4 x_5$ terms, the OLS regression is effectively specified
correctly, and the residual variance contracts proportionally to
the recovered signal. The same mechanism operates on the
quadratic DGP, where \epds\ reduces RMSE substantially from
\pds-\lasso's $0.231$, a reduction of nearly seventy percent.
The mechanism is what Proposition \ref{prop:erm} describes.
The true nuisance functions on DGPs 1 through 3 lie in $\Hcal$
at low complexity, the symbolic search has budget enough to
locate the Pareto elbow, and the recovered terms enter the
enriched \lasso\ step in a form that effectively spans the
truth.

On the linear-confounding design (DGP-5), \pds-\lasso\ holds at
bias $0.082$, indicating that the linear propensity score is
within the \lasso's representation capacity. The \pds\ structure
was designed for precisely this case, and the method recovers
the treatment effect with a small bias and conservative coverage.
\dml-NN exhibits notable bias of $0.238$, which we attribute to
the convergence-rate problem discussed in Section
\ref{sec:theory}: the neural network is over-parameterised
relative to the sample size, the implied sample-complexity
requirement is not met, and the residuals fail to absorb the
confounding. \dml-\lasso\ is slightly better at $0.115$,
reflecting the fact that the linear nuisance estimator is
appropriate for the linear propensity score in DGP-5.

The collapse of the standard methods on the nonlinear-confounding
designs (DGPs 6 and 7) is what the theory predicts. The
propensity score $m(\bm{x}) = 2 x_4 x_5$ lies outside the linear
span of $\bm{x}$, so the linear-\lasso\ hypothesis class does
not contain it, and the corresponding ERM bound under linearity
has an approximation error proportional to $\norm{m}^2$ that does
not vanish with $n$. For the neural-network nuisance estimator,
the hypothesis class does contain $m$ in principle, but the
sample-complexity calculation gives an implied $n$ on the order
of $10^6$ for the $(64, 32)$ MLP, so the rate-of-convergence
guarantee in \citet{chernozhukov2018} does not bind at the
sample sizes we study. \pds-\lasso, \dml-\lasso, and \dml-NN
exhibit bias between $1.20$ and $1.27$ on DGP-6, and between
$1.53$ and $1.59$ on DGP-7. These are statistically
indistinguishable from the naive-OLS bias of $1.24$ and $1.71$
respectively, indicating that none of the standard methods makes
any progress on the underlying confounding. \epds\ sidesteps both problems. Its hypothesis class $\Hcal$
contains the truth at low complexity on DGP-6 ($x_4 x_5$ has
complexity 3) and at moderate complexity on DGP-7 ($x_1^2 x_2$
adds complexity 5 for a total of 11 nodes), and the log-%
cardinality of $\Hcal$ is bounded, so the ERM rate is parametric. The \epds\ variants
reduce bias by more than ninety-five percent on DGP-6 and by
more than eighty-five percent on DGP-7, with corresponding
reductions in RMSE. On DGP-6, RMSE drops from $1.223$ for
\pds-\lasso\ to a substantially smaller value for \epds, a reduction of nearly
ninety percent.

On DGP-7 the two variants diverge in an informative way:
\epds\ shows a particular bias-variance trade-off, with bias slightly
higher RMSE ($0.713$ versus $0.677$). Decomposing
$\text{RMSE}^2 = \text{Bias}^2 + \text{Var}$ gives implied
variances of approximately $0.45$ for \epds\ and $0.48$ for
\epds. The method therefore reduces systematic error at the
cost of a small variance penalty, plausibly because the Du
screening occasionally retains or drops a variable that would
have moved through the symbolic-regression step under the
unfiltered procedure. We expect the variance penalty to shrink
with $n_{\text{rep}}$ since it is dominated by replication-level
noise in the screening step, but flag the divergence as one of
the finer comparisons that the production run should pin down.

\begin{table}[H]
\centering
\caption{Coverage of nominal $95\%$ confidence intervals at
$n=500$. Pilot results, $n_{\text{rep}} = 10$.}
\label{tab:coverage}
\small
\setlength{\tabcolsep}{4pt}
\begin{tabular}{lrrrrrrr}
    \toprule
    DGP
        & Full
        & \pds-\lasso
        & \dml-\lasso
        & \dml-RF
        & \dml-NN
        & \epds
        & \epds (CF) \\
    \midrule
    1 Linear
        & $1.000$ & $1.000$ & $1.000$ & {\small TBD} & $0.900$ & $1.000$ & {\small TBD} \\
    2 Quadratic
        & $1.000$ & $1.000$ & $1.000$ & {\small TBD} & $0.900$ & $1.000$ & {\small TBD} \\
    3 Interactions
        & $1.000$ & $1.000$ & $1.000$ & {\small TBD} & $0.900$ & $1.000$ & {\small TBD} \\
    4 Mixed NL
        & $1.000$ & $1.000$ & $1.000$ & {\small TBD} & $1.000$ & $0.800$ & {\small TBD} \\
    5 NL + linear conf.
        & $1.000$ & $1.000$ & $1.000$ & {\small TBD} & $0.900$ & $0.900$ & {\small TBD} \\
    6 NL + mild NL conf.
        & $0.000$ & $0.000$ & $0.000$ & {\small TBD} & $0.000$ & \hl{$1.000$} & {\small TBD} \\
    7 NL + severe NL conf.
        & $0.000$ & $0.000$ & $0.000$ & {\small TBD} & $0.000$ & \hl{$0.900$} & {\small TBD} \\
        8 Weak NL conf. (dummies)
        & {\small TBD} & {\small TBD} & {\small TBD} & {\small TBD} & {\small TBD} & {\small TBD} & {\small TBD} \\
    \bottomrule
\end{tabular}
\end{table}

The coverage results sharpen the picture. On DGPs 6 and 7, every
standard method has coverage of exactly zero. In all ten
replications on each DGP, the $95\%$ confidence interval misses
the true $\beta_0$. This is not a function of confidence-interval
width. The intervals are narrow, but they are centred on badly
biased point estimates. \epds\ achieve coverage of
$1.000$ on DGP-6 and $0.900$ on DGP-7. The single miss on DGP-7
reflects the residual bias from incomplete recovery of the cubic
confounder $2 x_1^2 x_2$. Two ingredients of the theoretical
story explain this. The cubic has complexity 5 in \pysr's
representation, near the upper end of what 40 iterations can
reliably locate, and the propensity-equation parsimony coefficient
$\rho_d = 0.01$ is calibrated to avoid spurious terms but
penalises the cubic at the margin. We expect higher recovery
rates and consequently higher coverage at the planned full
production replication count, and possibly higher still after
increasing the \pysr\ iteration budget.

On DGP-4, the \epds\ variants both have coverage of $0.800$, the
lowest value observed for these methods across the seven DGPs.
The reason is the operator-set mismatch for the $\log|x_3|$
term, which corresponds in the theory to the truth lying
strictly outside $\Hcal$: the Pareto elbow lies at the best
approximation rather than at the true expression, and the
resulting omission introduces a small bias that the $95\%$
confidence interval failed to absorb in two of ten replications.

\subsection{Trajectory from $n=30$ to $n=5000$}

The fixed-$n$ result raises a natural question. Does \dml's
failure on nonlinear confounding fix itself with more data, in
line with the asymptotic theory in \citet{chernozhukov2018}? Or
is the failure structural? To answer this we ran the full method
comparison on each DGP at ten sample sizes spanning two and a
half orders of magnitude, with three replications per setting.
The results are in Figure \ref{fig:trajectory}.

\begin{figure}[H]
    \centering
    \includegraphics[width=\textwidth]{figures/n_trajectory.png}
    \caption{Bias (left column) and RMSE on log scale (right
    column) versus sample size, one DGP per row. Lines show
    methods; shaded bands show $\pm 1$ Monte Carlo standard
    deviation across $n_{\text{rep}} = 3$ replications. Bands
    are wide and should be read as illustrative only at this
    pilot replication count; the production run will provide
    tight bands at $n_{\text{rep}} = 500$. Bottom two rows are
    the nonlinear-confounding designs.}
    \label{fig:trajectory}
\end{figure}

The answer is unambiguous. On the mild nonlinear-confounding
design (DGP-6, second from bottom), the bias of \dml-\lasso\ and
\dml-NN sits at approximately $1.1$ to $1.2$ from $n=30$ all the
way to $n=5000$. There is no visible improvement, and the
lines are essentially horizontal across the full $n$ range. The RMSE on
the same design hovers between $1.0$ and $1.3$ for the standard
methods across the entire range, on a log scale that would
otherwise reveal even modest improvements. The \epds\ variants
drop to zero bias by $n=200$ and to RMSE below $0.2$ by
$n=1000$. The difference is not a matter of degree, but a matter
of whether the method makes any progress at all.

The flat-bias pattern for \dml-NN is the empirical signature of
an unmet rate condition. The sample-complexity calculation in
the discussion following Proposition \ref{prop:erm} predicts a
required $n$ on the order of $10^6$ for the $(64, 32)$ MLP to
deliver $\epsilon = 0.1$ accuracy on a genuinely nonlinear
propensity score. Our trajectory caps at $n=5000$, three orders
of magnitude short of this threshold. The empirical observation
that bias does not decline with $n$ across the available range
is therefore not a small-sample artefact but a manifestation of
the rate condition not yet engaging.

On the severe nonlinear-confounding design (DGP-7), the same
pattern holds but with a slower convergence for \epds. The
standard methods remain flat at bias around $1.5$ across all
$n$, again with no visible improvement. \epds\
descend more gradually, reaching bias below $0.2$ by $n=1000$ on
\epds. This is consistent with
the genetic-programming convergence argument. The cubic term
$x_1^2 x_2$ has complexity 5 and is therefore higher up the
Pareto frontier than the interaction terms in DGP-6, so the
search budget binds more tightly and the convergence to the
correct expression is gradual rather than sharp. One anomaly
in the figure should be flagged. At $n=4000$ and
$n=5000$,
the bias of both \epds\ variants rises again on DGP-7. The pattern is
consistent across replications, which rules out an
issue, and is consistent with a small number of replications
hitting a \pysr\ failure mode at large $n$ where the search runs
out of iterations before converging on the cubic term. At higher
$n_{\text{rep}}$ in the planned production run we expect this
artefact to wash out, or alternatively to remain visible and
require an explanation. Until then we report it honestly rather
than smoothing it over.

The trajectory analysis also clarifies the linear DGP. \dml-NN's
bias on DGP-1 is small but consistently nonzero from $n=200$
onward (the red line in the top-left panel hovers slightly above
the bias-zero line through the full $n$ range), whilst other
methods are essentially unbiased. This is a hyperparameter-%
sensitivity story. A $(64, 32)$ MLP with default early-stopping
is not the right tool for a five-variable linear problem, and
the orthogonalisation step amplifies the resulting fitting
noise. The RMSE panel for the Linear DGP shows the same pattern:
\dml-NN's RMSE descent is shallower than \pds-\lasso's, and at
$n=5000$ it remains roughly twice as high. The result should be
read as evidence that \dml-NN is fragile to hyperparameter
choice, not as evidence against the \dml\ framework as such.

On the easy DGPs (Linear, Quadratic, Interactions), \epds\
achieves the lowest RMSE at essentially all $n$. The Interactions
panel is particularly striking. By $n=1000$, \epds\ has RMSE
around $5 \times 10^{-3}$, an order of magnitude below
\pds-\lasso\ and two orders of magnitude below \dml-NN. The
mechanism, as above, is that \pysr\ discovers the interaction
terms and the OLS on the enriched dictionary is effectively
correctly specified.

On the Mixed nonlinear design (DGP-4), the \epds\ variants do not
dominate. RMSE is comparable to \pds-\lasso\ across the entire
$n$ range, and at $n=5000$ \dml-NN is roughly tied with both
\epds\ variants. This is the design where the truth lies outside
$\Hcal$: \pysr\ recovers the cubic and interaction terms
reliably but recovers the log term in only half the cases, and
the resulting residual misspecification limits the bias
reduction. The story is honest: \epds\ dominates when the true
basis is discoverable within the available operator set, and
matches the standard methods when the operator set is mismatched.

\subsection{Term Recovery}

Figure \ref{fig:recovery_n} shows the fraction of replications
in which \pysr\ recovers each true nonlinear term, as a function
of $n$.

\begin{figure}[H]
    \centering
    \includegraphics[width=0.95\textwidth]{figures/recovery_vs_n.png}
    \caption{Recovery rates of true nonlinear terms versus sample
    size, one DGP per row. Solid lines are post-LASSO recovery;
    dashed lines are pre-LASSO (PySR discovery). The two coincide
    in most cases, indicating LASSO retains what PySR finds.}
    \label{fig:recovery_n}
\end{figure}

Three patterns emerge from the figure. First, on the
Quadratic and Interaction DGPs, recovery climbs steeply from
near zero at $n=30$ to near one by $n=200$ for the larger-%
coefficient term, and to between $0.5$ and $0.8$ by $n=500$ for
the smaller-coefficient term. The recovery story aligns directly
with the bias story. As \pysr's recovery rate of the active
terms increases, the residual confounding bias decreases. The
two panels for these DGPs are mirror images of the corresponding
bias trajectories, which is the kind of consistency check we
want from the recovery diagnostic. In the theoretical language,
recovery climbing to one is the GP search succeeding, and the
ERM bound in Proposition \ref{prop:erm} kicking in.

Second, the pre-LASSO and post-LASSO recovery rates coincide in
nearly every panel. This is reassuring, and it means the
enriched \lasso\ retains what \pysr\ discovers and does not drop
it on spurious-shrinkage grounds. The two-stage character of \epds,
which might in principle have suffered an interaction where
\lasso\ undoes \pysr's discoveries, does not exhibit that
interaction in practice. The pre and post lines in the figure
are visually indistinguishable in most panels, which we
interpret as evidence that the LASSO selection step on the
enriched dictionary behaves the way the theory suggests it
should.

Third, recovery of the $\log|x_3|$ term is persistently weak
across all DGPs where it appears, climbing only to around $0.5$
even at $n=5000$. The reason is the operator-set mismatch in
\pysr. The \texttt{log} operator computes $\log(x)$ for $x > 0$,
but our $x_3 \sim N(0, 1)$ takes negative values, so the term
must be discovered as $\log(x_3^2)$ rather than directly as
$\log|x_3|$. The parsimony penalty often rules out the more
complex form. This is the regime of the theory in which the
truth lies outside $\Hcal$, and the Pareto elbow corresponds to
the best approximation rather than to the truth. The fix is
either to add an absolute-value operator to \pysr\ or to draw
$x_3$ from a strictly positive distribution. We defer this to
the production run.

On the nonlinear-confounding DGPs (rows 5 and 6 of the figure),
recovery of the $x_4 x_5$ term reaches one across nearly the
full $n$ range. This is the smoking-gun result: \pysr\ reliably
discovers the nonlinear confounder, and the post-LASSO step
retains it, which is why the bias on these DGPs collapses to
zero for \epds. The cubic term $x_1^2 x_2$ on DGP-7 is harder,
with recovery rates in the range of $0.5$ to $0.75$ across the
upper $n$ values, which explains why the corresponding bias
trajectory for DGP-7 is slower to converge than for DGP-6 and
why the search-budget argument from Section \ref{sec:theory}
manifests visibly in the data.

\subsection{Visual Summary}

\begin{figure}[H]
    \centering
    \includegraphics[width=\textwidth]{figures/summary_panel.png}
    \caption{Bias, RMSE, coverage, and size across all DGPs and
    methods at $n=500$. The dramatic separation on DGPs 6 and 7
    between the \epds\ variants and the standard methods is
    visible across all four panels.}
    \label{fig:summary}
\end{figure}

% ============================================================
\section{Discussion}
\label{sec:discussion}

\subsection{When \epds\ Helps and When It Doesn't}

The simulation evidence and the theoretical story behind it
support a clean answer to the practical question of when to use
\epds. On settings with sparse linear nuisances, the truth lies
in the linear span and \pds-\lasso's hypothesis class contains
the truth, and \epds\ adds no flexibility that is needed and
imposes no measurable cost. On settings with discoverable nonlinear
structure but no nonlinear confounding, the truth lies in
$\Hcal$ at low to moderate complexity, the symbolic search
locates it, and \epds\ offers a modest RMSE improvement through
cleaner specification, with \epds\ achieving the
largest gains. On settings with nonlinear confounding, \epds\
is essentially the only method that works. The
standard alternatives operate in the wrong hypothesis class
(linear) or with a class whose sample-complexity requirements
do not bind at applied $n$ (neural networks).

The setting where \epds\ does not clearly dominate is DGP-4
(mixed nonlinear without confounding). RMSE is comparable to
\pds-\lasso, and coverage drops to $0.800$. The root cause is
the log-operator mismatch, which in the theoretical language
corresponds to the truth lying strictly outside $\Hcal$. With a
corrected operator set or a strictly positive distribution for
the variable inside the logarithm, the truth would re-enter
$\Hcal$ and we expect this gap to close. For applied work where
the suspected functional forms include absolute-value-style
logarithms, the operator set should be configured accordingly.
The broader lesson from the theoretical analysis is that the
success of \epds\ is contingent on the \pysr\ operator set
covering the truth, and this is the choice the analyst needs to
make thoughtfully.

\subsection{Why \dml\ Does Not Recover at Large $n$}

The trajectory result that \dml-NN bias stays above $1.0$ on the
nonlinear-confounding designs even at $n=5000$ deserves
unpacking, because it illustrates the sample-complexity argument
quantitatively. The convergence-rate condition in
\citet{chernozhukov2018} requires $\norm{\hat g - g}_2 \cdot
\norm{\hat m - m}_2 = o_p(n^{-1/2})$. For a multilayer
perceptron of architecture $(64, 32)$ on a $50$-dimensional
input, the implied parametric complexity is on the order of
$5000$ free parameters. The classical neural-network capacity
bound \citep{anthony2009neural} gives a sample-complexity
requirement scaling as $n = \Omega(P^2 \log P / \epsilon^2)$,
which at $P = 5000$ and $\epsilon = 0.1$ delivers a required $n$
on the order of $10^6$. Adding to this the fact that with five-%
fold cross-fitting each nuisance estimator sees only $4000$
observations during training, the effective sample size is even
smaller.

The practical implication is that \dml-NN's theoretical
guarantees do not bind in the moderate-$n$ regime that
characterises most applied microeconometric work. The procedure
can still be used, but its output should be cross-checked
against an estimator that does not rely on universal
approximation. \epds\ provides such a check. When \dml-NN and
\epds\ disagree by more than a standard error on data where
nonlinear confounding is plausible, the discrepancy should be
taken seriously, and a researcher should not default to the
\dml\ result on grounds of its more flexible nuisance estimator.

A reviewer may reasonably ask whether better neural-network
hyperparameters would rescue \dml-NN. We expect that a much
larger training budget, more aggressive regularisation, and a
network architecture explicitly tuned to the dimensionality of
the problem would improve performance. But the broader point
survives any particular hyperparameter choice: \dml-NN's success
depends on the analyst tuning the nuisance estimator well, which
itself requires either substantial domain knowledge or
substantial training time. \pds-\lasso\ and \epds, by contrast,
are essentially parameter-free for the user: $\lambda$ is set by
the theoretical formula in \citet{bch2014}, and the \pysr\
parsimony coefficient is the only meaningful choice. The trade-%
off is not just bias against variance, but also bias against
analyst effort.

\subsection{Why \epds\ Succeeds at Moderate $n$}

The theoretical analysis in Section \ref{sec:theory} can be
summarised in one sentence: \epds\ works because symbolic
regression searches a finite class with bounded log-cardinality,
and finite-class ERM enjoys parametric rates of convergence. The
relevant comparison is not whether \pysr\ can approximate any
function as well as a neural network can, but whether \pysr\
finds the structured truth faster than a neural network learns
it. For the structured class of functions we study, the answer
is yes by orders of magnitude, and the order-of-magnitude
calculation is set out explicitly in the comparison of $n =
O(\log |\Hcal|)$ with $n = \Omega(P^2 \log P)$.

The general principle is bias-variance. When the truth lies in
a small structured class, an estimator that imposes the class
as a prior dominates an estimator that treats the function as
an arbitrary element of a much larger class. At moderate sample
size, structural restriction wins. At very large sample size,
universal approximation eventually catches up, but our trajectory
analysis suggests that the crossover point lies well above
$n=5000$ for the DGPs we consider, and the sample-complexity
calculation above explains why this is so. This is a positive
result for applied work. Applied datasets typically have $n$ in
the hundreds or low thousands, and at those sample sizes
structural restriction is the right tool.

\subsection{Why Symbolic Methods Matter in Economics}

The case for preferring a symbolic-regression-augmented
procedure over neural-network \dml\ is not exclusively about
bias and coverage. It is also about the scientific virtues that
empirical economics has historically prioritised and that
black-box methods undermine. We discuss four.

\subsubsection{Falsifiability}

Popper's criterion for a scientific
claim is that it should be falsifiable. \pds\ produces claims
that meet this criterion by construction. The procedure names a
specific set of controls, and any subsequent researcher can
refit the regression with those controls and test whether the
estimated effect replicates. The symbolic-regression augmentation
preserves this property whilst extending the admissible nuisance
functions. The estimate $\betahat$ is now associated with a
specific set of original variables and a specific set of
discovered nonlinear terms, all of which can be inspected,
criticised, and rerun. \dml-NN, by contrast, produces an
estimate associated with a black-box $\hat{m}$ and $\hat{g}$
that cannot meaningfully be inspected by a second researcher.
The estimate is reproducible only in the weak sense that the
code can be rerun, not in the stronger sense that the underlying
claim can be evaluated and tested.

\subsubsection{Reproducibility}

Given the same data, the same operator
set, and the same parsimony coefficient, the symbolic-regression
step produces the same discovered terms. The result is
deterministic conditional on a small number of hyperparameter
choices that the analyst writes down explicitly. \dml-NN, by
contrast, depends on initialisation, batch ordering, GPU non-%
determinism, and a host of implementation details that typically
go untracked. The reproducibility concerns in applied work are
partly a reflection of black-box methods whose outputs depend on
factors the analyst does not document.

\subsubsection{Narrative}

An applied paper in labour economics,
development, or finance using \pds\ can include a sentence like
``the procedure selected baseline income, education, household
size, and an interaction between age and labour-force
participation as relevant controls.'' This sentence is readable
by the policy audience, by referees, by future meta-analysts,
and by anyone trying to understand why the treatment effect is
what it is. The same sentence cannot be written about \dml-NN.
The narrative role of named controls is not merely cosmetic. It
is the mechanism by which an empirical claim becomes part of
cumulative scientific knowledge rather than a one-off output of
a pipeline.

\subsubsection{Sensitivity analysis}

When the analyst can see the
controls, they can perturb them. They can ask what happens to
$\betahat$ if a particular discovered term is dropped, or what
happens if a control the procedure did not select is added back.
The sensitivity tools of \citet{cinelli2020sensitivity} and
\citet{conley2012plausibly} all rely on being able to vary the
assumed structure of the nuisance functions. These tools cannot
be applied to \dml-NN in any meaningful way, because the
nuisance functions are not specifiable objects. A method that
admits sensitivity analysis is a method that can be defended
against scepticism, and in applied work that distinction
matters.

Falsifiability and narrative are particularly pressing in
economics because economic claims typically need to support
both academic peer review and policy deliberation. A treatment
effect that cannot be defended in plain language to a
non-technical audience is rarely useful to the people who would
act on it, however technically valid the estimate.
Symbolic-regression augmentation restores the named-control
property that \pds\ preserves and \dml-NN gives up, and does so
without sacrificing the representational power that motivated
introducing flexible learners in the first place.

\subsection{Where Symbolic Regression Falls Short}

Symbolic regression is not a universal tool. Its failure modes
are part of taking the framework seriously, and four of them
are worth setting out.

\subsubsection{Mismatched operator sets}

Our DGP-4 result is the
canonical illustration. When the true relationship involves
transformations outside the configured operator set, the
symbolic search cannot recover the truth even with arbitrary
data. The Pareto elbow then sits at the best approximation
within the class, and the resulting bias does not vanish with
$n$. This is the regime where Proposition \ref{prop:erm}'s
approximation-error term binds. In practice, the analyst's
choice of operator set is a substantive specification decision,
comparable to choosing which controls to include in a linear
regression, and the procedure inherits no protection against
getting that choice wrong.

\subsubsection{Genuinely non-symbolic truth}

Symbolic regression assumes the truth is the kind of function a domain expert might
write down, namely a short combination of recognisable
operators. When the truth is instead the kind of function a
flexible learner happens to produce, an arbitrary smooth
surface with no compact symbolic representation, the search
will return either a degenerate constant or a high-complexity
expression that overfits noise. The \dml\ programme with a
well-tuned learner is designed for exactly this case, and we
should not expect symbolic regression to compete in regimes
where the structural assumption is false.

\subsubsection{High dimensionality}

\pysr's runtime and the spurious-term problem both worsen with
the number of input variables. We tested up to $p=50$, where a
single \pysr\ fit takes approximately twenty seconds and
discovers true terms reliably under the parsimony budget reported
in Section \ref{sec:method}. At very large $p$ the symbolic-%
regression step is computationally prohibitive, and an upstream
variable-screening step (for example a Stein-based selector along
the lines of \citealp{duetal2025}, or a sparsity-inducing
neural network) would be needed to keep the search tractable.
We would not extrapolate beyond $p \approx 100$ without a more
careful analysis of computational and statistical scaling, and
we discuss upstream screening as future work in Section
\ref{sec:future}.

\subsubsection{Low signal-to-noise}

Symbolic regression is a
model-discovery tool, and like any such tool it depends on a
clean signal. At low SNR the Pareto frontier loses its sharp
elbow and the parsimony coefficient becomes consequential
beyond what theoretical guidance can fix. The DGPs we
construct have favourable SNR by design. In applied work it
will often be worse, and an analyst should expect the
procedure to discover fewer terms and recover less of the
truth as a result.

These failure modes are not unique to our particular
implementation. They are properties of symbolic regression as a
method, and they limit the scope of the framework
correspondingly. The procedure works where the truth is
structured, the operator set is well-chosen, the dimensionality
is moderate, and the signal is reasonably clean. Where any of
these conditions breaks, the analyst should consider whether
\dml\ with a flexible learner, or \pds-\lasso\ if the linearity
assumption is defensible, is a better tool for the problem at
hand.

\subsection{Limitations}

Six limitations of the current implementation are worth
flagging.

\subsubsection{No sample splitting}

The formal \citet{bch2014} theorem
assumes a fixed basis dictionary. We estimate the basis from the
same data we use for the final regression. The simulation
evidence on coverage is encouraging but does not constitute a
proof of inferential validity. A sample-split variant, where
\pysr\ discovers the basis on half the data and the
\pds-\lasso\ step runs on the other half, would restore formal
validity at the cost of statistical efficiency. We defer the
implementation and the analysis of the resulting efficiency loss
to future work, but flag this as the principal outstanding
inferential question.

\subsubsection{Operator-set mismatch for the logarithm}

As discussed, \pysr's \texttt{log} operator does not naturally discover
$\log|x_i|$ when $x_i$ takes both positive and negative values.
This is the single largest source of imperfect coverage in our
results (DGP-4, coverage of $0.800$). The truth lies outside
$\Hcal$ for our operator-set configuration, which is the
misspecified regime of Remark \ref{rem:misspec}. A clean fix is straightforward and will be implemented before
the production run, either by adding an absolute-value operator
to the \pysr\ operator set, or by drawing covariates that
appear inside logarithms from strictly positive distributions
in the DGPs.

\subsubsection{Small pilot replication count}

The fixed-$n$ tables use
$n_{\text{rep}} = 10$ and the trajectory analysis uses
$n_{\text{rep}} = 3$. The Monte Carlo standard errors are wider
than we would prefer, particularly for the trajectory. The
directionally large effects we emphasise are robust to the
replication count. Bias of $1.2$ versus $0.0$, or coverage of
$0.0$ versus $0.9$, is not going to flip sign at higher
replication. The finer comparisons, such as the precise
convergence point of the trajectory, should be re-examined at
$n_{\text{rep}} = 500$.

\subsubsection{No applied analysis}

The dissertation is simulation-based. An applied analysis on real economic data with
plausibly nonlinear confounding has not been carried out. The
\citet{bch2014} eminent-domain application is a natural
candidate, where the linear-controls specification is
plausibly missing nonlinear judge-characteristic or
circuit-court-composition interactions. Without an applied
analysis, the framework's practical relevance is supported by
simulation evidence only, and a reader who weighs real-data
performance heavily in their assessment of a method should
treat the contribution accordingly.

\subsubsection{Computational cost}

Each \pysr\ fit takes
approximately twenty seconds at our settings. For a single
applied analysis this is trivial, but the full $n_{\text{rep}}
= 500$ production run requires high-performance computing
infrastructure to be tractable on the time scales relevant to
dissertation work. The cost compounds for any analysis
requiring repeated fits, including bootstrap inference,
sensitivity analysis under varying operator sets, and
cross-fitting. Reducing the cost is partly a matter of \pysr\
hyperparameter tuning and partly a matter of advances in the
symbolic-regression literature we discuss below.

\subsection{Practical Guidance}

For applied work, we offer three recommendations. When the
suspected confounders are believed to enter the outcome and
propensity equations approximately linearly, \pds-\lasso\
remains the right tool, since it is faster, simpler, and
theoretically well-understood. When nonlinear functional forms are suspected
but confounding is unlikely, both \pds-\lasso\ and \epds\ are
reasonable, and running both and comparing the discovered
nonlinear terms is itself informative as a specification check.
When confounding is suspected to be nonlinear, particularly when
the treatment is plausibly a function of interactions or higher-%
order terms in the covariates, \epds\ is the preferred procedure.
The computational cost is moderate
(approximately twenty seconds per run at $n=500, p=50$), and the
gain in bias and coverage is large enough to justify it.

A useful diagnostic in applied work is to inspect the discovered
\pysr\ expressions. If \pysr\ returns expressions involving
products, quadratics, or logarithms of variables that
substantively make economic sense as confounders, the analyst
gains confidence in the result. If \pysr\ returns expressions
involving variables that have no plausible substantive role, the
analyst should investigate whether the discovered terms are
genuine signal or artefacts of noise. The interpretability of the
\pysr\ output is part of the value of the method, in that it
gives the analyst a window into the data-generating process that
the black-box \dml-NN procedure does not.

\subsection{A Framework Beyond Post-Double Selection}

The contribution of this dissertation generalises beyond the
\pds\ pipeline. The underlying idea, that nuisance functions in
causal inference admit a structured-versus-flexible choice and
that structural restriction can dominate when the truth genuinely
lies in the structured class, applies wherever a causal procedure
requires functional-form decisions on auxiliary objects.

Three extensions follow naturally. First, in propensity-score
matching and inverse-probability weighting, the propensity model
is itself a nuisance function whose functional form can be
misspecified. A symbolic-regression-augmented propensity step
could recover nonlinear treatment-assignment structure that a
logit-with-main-effects baseline misses, and the downstream
weighting would then operate on a propensity estimate that
respects the assignment mechanism rather than a linear
approximation to it. The interpretability of the discovered
expressions is particularly valuable here. A policy researcher
who can name the determinants of treatment assignment is in a
better position to defend the resulting estimate than one who
has produced a black-box propensity score.

Second, in difference-in-differences and event-study designs,
the conditional outcome trajectory under no treatment is a
nuisance function. The well-known concern about parallel
trends, that pre-trend differences may reflect nonlinear
baseline heterogeneity rather than evidence of common dynamics,
is partly a functional-form concern. Replacing a linear
pre-trend control specification with a symbolic-regression-%
recovered functional form would address this concern directly,
at the cost of accepting the operator-set assumption that
follows.

Third, in instrumental variables with many candidate
instruments, the first-stage relationship is a nuisance
function. The IV-\lasso\ procedure of \citet{bch2014} and the
broader instrument-selection literature face the same
approximate-sparsity assumption that \pds\ does, and the same
symbolic-regression augmentation would address it. The
structural-restriction advantage we document for \pds\ should
transfer cleanly to the IV setting.

The common thread is that any causal procedure that selects,
partials out, or otherwise depends on an estimated nuisance
function can in principle be augmented by a symbolic-regression
discovery step, provided the analyst is willing to specify an
operator set and accept its coverage as a substantive
assumption. None of these extensions are trivial. Each involves
its own inferential considerations and operator-set choices.
But the methodological recipe is the same as the one we have
implemented and evaluated in the \pds\ setting. The framework
is more general than any particular instantiation of it, and
the \pds\ application should be read as a worked example rather
than the contribution itself.

% ============================================================
\section{Conclusion}
\label{sec:conclusion}

We have proposed \epds, a symbolic-regression-augmented
extension of
the \citet{bch2014} \pds-\lasso\ framework that uses symbolic
regression to discover nonlinear basis terms directly from the
data. The method preserves the principal virtue of \pds, namely
a named and interpretable set of selected confounders, whilst
extending the representational capacity to nonlinear nuisance
functions. We position the framework as a single method rather than
the Stein-based feature selector of \citet{duetal2025} to reduce
the dimensionality of the symbolic search.

The theoretical reason \epds\ works is that symbolic regression
performs empirical risk minimisation over a finite hypothesis
class with bounded log-cardinality, so classical concentration
arguments deliver a parametric rate of convergence when the
truth lies in the class. The corresponding sample-complexity
requirement for a neural-network nuisance estimator scales
quadratically in the parameter count, putting the required $n$
beyond the reach of applied work. This is the formal
counterpart to the empirical observation that, on nonlinear-%
confounding designs, the \dml-NN bias does not decline with $n$
across the range we examine. \epds\ avoids this by restricting
the search to a small structured class where the rate condition
binds.

Across seven data-generating processes, with particular
attention to two novel designs involving nonlinear confounding,
the proposed methods deliver substantial improvements in bias,
RMSE, and coverage over \pds-\lasso, \dml-\lasso, and \dml-NN,
in the regime where the standard methods fail. The trajectory
analysis from $n=30$ to $n=5000$ shows that \dml's failure on
nonlinear confounding is not a small-sample problem that
resolves with more data, but a manifestation of the convergence-%
rate condition not yet engaging at the sample sizes typical of
applied work. \epds\ converges to unbiasedness by $n=200$ on
the mild nonlinear-confounding design and by approximately
$n=1000$ to $n=2000$ on the severe nonlinear-confounding design.

\subsection{Future Directions}

The work points to several directions for further development.
We group them by what kind of effort each requires:
theoretical, practical, frontier, and applied.

\subsubsection{Theoretical questions}

The most pressing is a formal
inference theorem for the symbolic-regression-augmented
dictionary. A Belloni-style result accommodating data-estimated
bases, perhaps via the uniform post-selection inference
machinery of \citet{tibshirani2018uniform} or via a
sample-splitting argument, would put the inferential validity
of the procedure on rigorous footing. Adjacent to this is the
question of optimal complexity-penalty selection. The parsimony
coefficient $\rho$ in our procedure is set heuristically, and a
principled BIC-style derivation tied to sample size would close
a clear gap. Finally, the rate calculation in Proposition
\ref{prop:erm} assumes the truth lies in $\Hcal$. Extending the
analysis to the misspecified case, where the truth is close to
but not in the class, would clarify how much approximation
error the procedure can tolerate before its advantage over
\dml\ disappears.

\subsubsection{Practical challenges}

The operator-set choice is
currently a researcher degree of freedom, and getting it wrong
costs accuracy. Adaptive operator selection, perhaps starting
with a small set and expanding when residual structure remains
in the data, would reduce this burden. Computational cost is
the other practical bottleneck. Each \pysr\ fit takes order
twenty seconds at our settings, which is fine for a single
analysis but limiting for any simulation or bootstrap requiring
repeated fits. Integration with cross-fitting, which would
address some of the inferential concerns above, compounds the
computational cost further. The procedure has not been tested
at $p$ much above fifty, and scaling to genuinely
high-dimensional settings will require either tighter upstream
screening or more efficient symbolic search.

\subsubsection{Frontier directions in symbolic regression}

The symbolic-regression literature is itself advancing rapidly,
and several developments are directly relevant to the framework.
Neural-guided genetic programming, in the spirit of
\citet{cranmer2020}, uses a graph neural network to propose
mutations and achieves substantially faster convergence than
vanilla genetic programming. Reinforcement-learning-based
symbolic regression treats expression-tree construction as a
sequential decision problem, with the agent learning to
construct expressions that fit data efficiently.
Transformer-based approaches train large models to predict
expressions directly from numerical data, sidestepping
evolutionary search entirely. Integrating any of these into
the symbolic-regression step of our pipeline would not change
the inferential structure of the downstream \pds\ procedure but
could dramatically reduce the computational cost and improve
recovery rates on high-complexity targets. The frontier is
moving quickly, and a pipeline that treats the symbolic
regressor as a modular component will inherit those gains as
they arrive.

\subsubsection{Applied next steps}

Beyond these methodological
directions, an applied analysis is the most natural next step.
We have a specific candidate in the \citet{bch2014}
eminent-domain application, where the linear-controls
specification is plausibly missing nonlinear interactions
between judge characteristics, circuit-court composition, and
regional economic conditions. A replication exercise that
re-runs the original analysis with the symbolic-regression
augmentation would demonstrate the framework's practical
relevance, and if the discovered nonlinear terms make
substantive economic sense, would constitute a stronger
empirical case than simulations alone can make. A second
direction is the extension to heterogeneous treatment effects.
Using symbolic regression to recover the functional form of
the treatment-effect surface itself, rather than the nuisance
functions of the controls, would connect this work to the
heterogeneous-effects literature represented by
\citet{wager2018causal} and \citet{athey2016recursive}, and
would broaden the applicability of the framework substantially.

\subsection{The Broader Methodological Message}

The contribution of this dissertation extends beyond the
specific bias-coverage comparison. The dominant view in
applied causal inference is that flexible machine learning
should be used directly for nuisance estimation, with
theoretical guarantees flowing through to the resulting causal
estimate. We have argued for a different arrangement. Flexible
tools, including symbolic regression itself, are best used to
discover the structure in the data. That structure, represented
in a form simple enough to be inspected and defended, then
carries the inference. The pattern has been productive in the
physical sciences, where flexible learners are used as
discovery tools and symbolic expressions become the final
scientific output. The same logic applies to causal
econometrics, where the simple representation enables
falsifiability, reproducibility, narrative, and sensitivity to
assumed specification, all virtues that the dominant practice
gives up too quickly. The complex tool finds the answer; the
simple tool delivers it. The specific procedure we evaluate is
one instantiation of this principle, applied to the
post-double-selection setting. The principle itself is more
general, and the framework should extend naturally to other
causal procedures where nuisance functions need to be estimated
from data.

% ============================================================
\newpage
\bibliography{references}

\end{document}
